\clearpage
\fancyhead[L]{\small SECOND COMET OF THE YEAR 1805.}
\hypertarget{native-280-lower}{}
\begin{center}
SECOND COMET OF THE YEAR 1805.

\rule{0.6\linewidth}{0.4pt}

Monatliche Correspondenz. Vol. XIV. Pp. 181–186. August 1806.

\rule{0.6\linewidth}{0.4pt}
\end{center}

\begin{center}
\textsc{Gauss to von Zach.}
\end{center}

\begin{flushright}
Brunswick, 8 July 1806.
\end{flushright}

\noindent --- --- --- Since I wrote to you about the comet of 1805, the only further work I have done has been to compare the observations by \textsc{Bouvard} that you were kind enough to send me with both my parabolic elements (M.\,C., July issue, 1806, \href{https://www.digitale-sammlungen.de/en/view/bsb10538606?page=87}{p.~77}) [\hyperlink{native-276}{p.~271} of this volume] and my elliptical elements (ibid., \href{https://www.digitale-sammlungen.de/en/view/bsb10538606?page=89}{p.~79}) [\hyperlink{native-277}{p.~272} of this volume]. Here is the comparison: [The difference is R---B.]

\begin{center}\begingroup
\fontsize{8}{10}\selectfont
\setlength{\tabcolsep}{1pt}\renewcommand{\arraystretch}{1.12}
\begin{tabular}{l||c|c|c|c||c|c|c|c}
\shortstack{Time of obs.} & \multicolumn{2}{c|}{\shortstack{Calculated in\\the parabola}} & \multicolumn{2}{c||}{\shortstack{Difference in}} & \multicolumn{2}{c|}{\shortstack{Calculated in\\the ellipse}} & \multicolumn{2}{c}{\shortstack{Difference in}} \\
Paris 1805. & Longitude & Latitude & Longitude & Latitude & Longitude & Latitude & Longitude & Latitude \\
\hline
Nov. 16.\,45203 & $29^\circ\,21'\,30''$ & $30^\circ\,2'\,32''$ & $-2'\,7''$ & $-8''$ & $29^\circ\,22'\,48''$ & $30^\circ\,1'\,18''$ & $-0'\,49''$ & $-1'\,22''$ \\
17.\,36386 & $28\,51\,5$ & $29\,49\,40$ & $-0\,18$ & $+48$ & $28\,51\,55$ & $29\,48\,47$ & $+0\,32$ & $-0\,4$ \\
18.\,39927 & $28\,14\,28$ & $29\,32\,42$ & $+0\,10$ & $-2$ & $28\,14\,43$ & $29\,32\,8$ & $+0\,25$ & $-0\,36$ \\
23.\,32241 & $24\,41\,33$ & $27\,25\,5$ & $+0\,29$ & $-30$ & $24\,39\,54$ & $27\,25\,31$ & $-1\,10$ & $-0\,4$ \\
30.\,51095 & $15\,43\,40$ & $19\,25\,32$ & $+4\,0$ & $+4$ & $15\,42\,3$ & $19\,26\,26$ & $+2\,23$ & $+0\,58$ \\
Dec. 5.\,29581 & $2\,7\,19$ & $3\,20\,5$ & $+0\,8$ & $-40$ & $2\,7\,20$ & $3\,21\,26$ & $+0\,9$ & $+0\,40$ \\
\end{tabular}\endgroup\end{center}

\par\indent This comparison shows that the ellipse does not agree quite as well with these observations as it does with those from Marseille. At the same time, it seems to show 1) that the Marseille observations are better, and 2) that the parabola gives larger differences than the ellipse for the other observations as well. I must add that the differences given here for the ellipse are not really the true ones and still need considerable changes. The \emph{calculated} longitudes and latitudes here are free of parallax and aberration. I compared them with the observed ones after they had been cleared of parallax and aberration (as you sent them to me). But that reduction, which had already been made, depends on the distances of the comet from the Earth. The reduction due to parallax is inversely proportional to the distance, whereas the reduction due to aberration is directly proportional to it. The distances, however, turn out quite differently in the parabola and the ellipse. For example:

\begin{center}\begingroup
\fontsize{9}{11}\selectfont
\setlength{\tabcolsep}{2pt}\renewcommand{\arraystretch}{1.12}
\begin{tabular}{rlll}
Distance in the parabola: & Nov. 16 $=0.2489$, & Nov. 23 $=0.1761$, & Dec. 5 $=0.0638$ \\
in the ellipse: & $=0.1794$ & $=0.1269$ & $=0.0460$ \\
\end{tabular}\endgroup\end{center}
\SourceLeaf{281}%
\par\indent The reduction for parallax, which is particularly considerable here, is thus greater for the ellipse than for the parabola in about the ratio 7:5, while the reduction for aberration is smaller in the same ratio. The observations do not give these two reductions separately. Only their sum can be inferred from the figures, and even that only for the last four observations. The reductions would therefore have to be calculated afresh for the ellipse, which I did not take the trouble to do. For the parabola, however, the reductions will come out as \textsc{Legendre} assumed, since my parabolic elements differ so little from his. One can nevertheless see even from a rough estimate that the observation of 5 December, the only one that seems to be a meridian observation, would agree much better still with the elliptical elements after a more exact reduction. If I should yet get to see \textsc{Bouvard}’s original right ascensions and declinations, I will compare them with my elements in the same way as I did with \textsc{Thulis}’s. For the present I have set these calculations aside, since nothing more definite can be said about the comet until observations come to light that extend the time span further.

\par\indent I have not yet seen the work by \textsc{Legendre} that you mention in this connection. I deliberately made no effort to obtain it, so that in working on my method I would stay entirely within the chain of my own ideas. A few words that \textsc{de la Lande} drops in the latest \emph{Histoire de l'Astronomie}, 1805, \emph{méthode des moindres quarrés}, lead me to suspect that a principle I have already been using for twelve years in various calculations, and that I will also use in my book, though it is not essential to my method, --- that this principle has also been used by \textsc{Legendre}. In any case there can be few points of contact, since our works have very different objects. Determining a comet’s orbit, where one assumes a parabola, and determining a planet’s orbit, where one is to set aside all hypotheses, require quite different treatment. My book is devoted \emph{primarily} to the latter problem, though it will also contain much of my own work on the former.

\par\indent I am also extremely glad that I did not publish as early as 1802 the method I had used to calculate the orbits of Ceres and Pallas, despite the many requests I received. Since then I have kept working to improve the method itself, especially last winter, and its present form bears hardly any resemblance to its original form. To give an idea of this, I will mention just \emph{one} point. The problem is so involved that by its very nature the very first approximation must make some assumptions that are only approximately correct (for instance, in Dr.~\textsc{Olbers}’s method, the chords for the Earth and the comet are assumed to be divided by the intermediate radius vectors in the ratio of the time intervals). Assumptions of this kind therefore also underlie the first approximation in my method, and they depart less from the truth the shorter the time intervals are. The observations on which the first approximation is based must therefore not be taken too far apart, since otherwise the merely approximate assumption would leave the first calculation too far from the truth and require too many laborious repetitions of the methods of correction. Still, with my method as it was in 1802, I could use observations of Pallas 27 days apart from the outset. I would not have cared to go much further, however. My method is now such that recently, when I wanted to use the observations of \textsc{Oriani} from 1805 that you were kind enough to send me for an
\SourceLeaf{282}%
\emph{example} for my book, and therefore had to treat them as though I knew nothing at all about the orbit of Pallas, I was able at once and with the greatest success to take the outermost observations, 71 days apart, as my starting point. There is no doubt at all that I could have gone considerably further.

